
% Do not change these lines
\renewcommand{\abstractname}{\textbf{{\Large ABSTRACT}}}
\addcontentsline{toc}{chapter}{\textbf{\normalsize{\emph{ABSTRACT}}}}

\begin{abstract}\thispagestyle{plain}
Bangla users do not always meet language models through the script assumed by
standard benchmarks. The same Bengali question may be typed in native Bengali
script, in English, or in Latin-script Banglish. Recent Bengali benchmarks
increasingly cover knowledge, education, culture, inference, and social
interaction, but they rarely isolate this script choice while holding the
underlying task and answer fixed. This thesis studies whether orthography
itself changes large language model behavior. We construct controlled Bangla,
reviewed Banglish, and English variants of curriculum-style QA and math items
drawn from BEnQA and BanglaMATH, and evaluate models under a paired item-level
protocol.

On a 200-item validation slice, compact Qwen models show a consistent
reviewed-Banglish deficit. Qwen2.5-3B scores 54/200 in Bangla but 41/200 in
reviewed Banglish, Qwen2.5-7B 8-bit scores 65/200 versus 47/200, and Qwen3-4B
scores 80/200 versus 49/200. The gap is not explained by token count alone,
since reviewed Banglish is token-cheaper than native Bangla for the audited
Qwen tokenizers. A frozen v5 review of high-impact Banglish rows changes scores
by at most two items, improving benchmark quality without removing the
script-conditioned weakness.

Hosted-model and scale checks clarify the boundary. GPT-5.5 low nearly closes
the validation-200 gap under secondary scoring, while Claude Sonnet 4.6,
DeepSeek V4 Flash, and Groq-hosted Llama 3.3 70B preserve reviewed-Banglish
deficits under the same protocol. A 1,000-row human review audit extends the
BEnQA evidence base; 974 accepted or edited rows form a gold extension, giving
1,174 usable human-reviewed paired items when combined with the 200-item core.
On this extension, Qwen2.5-3B, Groq-hosted Llama 3.3 70B, and DeepSeek V4 Flash
all reproduce the English greater-than Bangla greater-than reviewed Banglish
ordering. Mitigation remains model-dependent: self-normalization helps one
Qwen model and hurts another, and generated alternate-script views require
preservation gates before they can support a deployable routing claim. Overall,
the thesis shows that Latin-script Banglish is an undermeasured orthographic
robustness challenge for Bangla LLM use. It is a real access path for Bangla
users, and its reliability must be measured directly rather than assumed from
native-script Bangla or English performance.
\end{abstract}

\endinput
